\honest{Tsuyoku Naritai! (I Want To Become Stronger)}{Eliezer Yudkowsky}{2007}

Orthodox Judaism has this saying: ``The previous generation is to the next one as angels are to men; the next generation is to the previous one as donkeys are to men.'' Since all its law came from Sinai, knowledge can only be lost in transmission, and modern rabbis cannot overrule ancient ones. A worm found in an apple may be eaten, because the ancient rabbis thought it grew from the apple, and no modern rabbi may say they knew nothing of biology. As a schoolboy I reasoned that Torah loses knowledge every generation and science gains it, so sooner or later science must surpass Torah. What matters most is that there be progress.

The Japanese phrase of the title means ``I want to become stronger.'' It expresses a feeling ``embodied more intensely in Japanese works than in any Western literature I've read.'' You say it when you resolve to become a professional Go player, after losing an important match, after winning one while still short of the top rank, or after becoming the greatest player ever while still thinking you can do better.

On Yom Kippur, Orthodox Jews recite a list of sins, the Ashamnu, and strike their chests at each word. You cannot skip ``we have stolen'' because you did not steal. The day is about ``confessing sins---not avoiding sins so that you have less to confess,'' and the list does not end with a promise to do better next year. Rationalists do the same when they say ``We are all biased, we are all irrational.'' My answer: ``Fine. Now tell me how you plan to become less biased, less irrational,'' and so on.

\nb{This is good advice, and it is the tradition's own. Maimonides, in his code of Jewish law, says that repentance means abandoning your sins, ``resolving in his heart, never to commit them again.'' Confessing without that resolve, he says, is like immersing in a ritual bath while holding a dead lizard. And the Talmud records a Yom Kippur confession that asks, ``May it be Your will that I may sin no more.'' The post presents its point as a correction of Judaism.}

Then a joke. On Yom Kippur the rabbi cries, ``God, I am nothing before you!'' So does the cantor. When the janitor does the same, the rabbi nudges the cantor: ``Look who thinks he's nothing.'' Do not take pride in confessing your flaws. Just as pride in your ignorance makes you slow to give it up when evidence comes, pride in your self-awareness makes you slow to give up your flaws. The time to rejoice is when you have a little less to confess. If we take pride in confessing, then when someone brings a plan for correcting a bias, ``We will shake our heads sadly and say, `You must not be very self-aware.'\,''

``Never confess to me that you are just as flawed as I am unless you can tell me what you plan to do about it.'' You will still have plenty of flaws afterwards. What matters is to do better, to take one more step forward.
